\NeedsTeXFormat{LaTeX2e}

\documentclass[portrait]{a0poster}

\usepackage{epsfig}
\usepackage{color}
\usepackage{transparent}
\usepackage{wallpaper}
\usepackage{amsmath}
\usepackage{amssymb}
\usepackage{calc}
\usepackage{bm}
\usepackage[utf8]{inputenc}
\usepackage{textcomp}
\usepackage{mathptmx}
\usepackage[scaled=.90]{helvet}
\usepackage{courier}
\usepackage{listings}

\hyphenation{re-so-na-tors na-no-tubes}

\pagestyle{empty}
\TileWallPaper{10.234cm}{11.594cm}{images/dokutitle}

\setlength{\parskip}{5pt plus 2pt minus 1pt}
\setlength{\parindent}{0pt}

\frenchspacing
\sloppy

\renewcommand{\familydefault}{\sfdefault}

% some color definitions
\definecolor{darkblue}{rgb}{0, 0, 0.5}
\definecolor{lightblue}{rgb}{0, 0, 0.8}
\definecolor{lightgray}{gray}{0.75}

\newcommand{\blue}{\color[rgb]{0.2,0.2,1}}
\newcommand{\lightblue}{\color{lightblue}}
\newcommand{\red}{\color{red}}
\newcommand{\darkgreen}{\color[rgb]{0,0.65,0}}

% Delft blauw 
\definecolor{background}{rgb}{1,1,1}
\definecolor{text}{rgb}{0, 0, 0.3}
\definecolor{heading}{rgb}{0, 0, 0.8}

\definecolor{invheading}{rgb}{0.8,0.8,1}

\pagecolor{black}

\lstset{
  basicstyle=\ttfamily,
  identifierstyle=\color[rgb]{0,0,0.5},
  stringstyle=\color[rgb]{0,0.5,0}\textit,
  commentstyle=\color{red}\textbf
}




% large heading in text
\newcommand{\heading}[1]{
  {\color{heading}\boldmath\textbf{\huge #1}}\\[\medskipamount]
}

% small heading in text
\newcommand{\smallheading}[1]{
  {\color{heading}\textbf{\Large #1}}\\
}

% very small heading in text
\newcommand{\verysmallheading}[1]{
  {\color{heading}\textbf{\large #1}}\\
}

% white bar across the column
\newcommand{\whitebar}{%
\vspace*{0.5cm}
\hspace*{-1.5cm}
\fcolorbox{white}{white}{\rule{30cm}{-0.7cm}}

\vspace*{1cm}
}

% white background for plots etc.
\newcommand{\whitebg}[1]{
 {
  \setlength{\fboxsep}{5mm}
  \setlength{\fboxrule}{1pt}
  \fcolorbox{black}{white}{\color{black}{#1}}
 }
}

% place a figure
\newcommand{\placefigure}[1]{ 
 {
  \vspace{-0.4cm}
  \begin{center}
  \hspace*{-1cm}\begin{minipage}{0.96\textwidth}
   \whitebg{
     \epsfig{#1,width=0.96\textwidth}
   }
   \vspace*{-1cm}
   {\\}
  \end{minipage}
  \end{center} 
  \vspace{-1cm}
 }
}

% place a listing
\newcommand{\placelisting}[1]{ 
 {
  \vspace{-0.2cm}
  \begin{center}
  \hspace*{-1cm}\begin{minipage}{0.96\textwidth}
   \footnotesize \renewcommand{\baselinestretch}{0.85}
   \whitebg{
   \parbox{0.97\textwidth}{
    \lstinputlisting[language=Perl]{#1}
   }
   }
   \renewcommand{\baselinestretch}{1}
   \vspace*{-1cm}
   {\\}
  \end{minipage}
  \end{center} 
  \vspace{-1cm}
 }
}

% place a listing
\newcommand{\placesmalllisting}[1]{ 
 {
  \vspace{-0.2cm}
  \begin{center}
  \hspace*{-1cm}\begin{minipage}{0.96\textwidth}
   \scriptsize \renewcommand{\baselinestretch}{0.8}
   \whitebg{
   \parbox{0.97\textwidth}{
    \lstinputlisting[language=Perl]{#1}
   }
   }
   \renewcommand{\baselinestretch}{1}
   \vspace*{-1cm}
   {\\}
  \end{minipage}
  \end{center} 
  \vspace{-1cm}
 }
}

% place a text
\newcommand{\placetext}[1]{ 
 {
  \vspace{-0.2cm}
  \begin{center}
  \hspace*{-1cm}\begin{minipage}{0.96\textwidth}
   \footnotesize \renewcommand{\baselinestretch}{0.85}
   \whitebg{
   \parbox{0.97\textwidth}{
    \lstinputlisting{#1}
   }
   }
   \renewcommand{\baselinestretch}{1}
   \vspace*{-1cm}
   {\\}
  \end{minipage}
  \end{center} 
 }
}




\begin{document}

\color{text}
\Large
%
\newlength{\seplength}
\newlength{\headerheight}
\newlength{\columnheight}
\newlength{\columnheighta}
\newlength{\columnheightb}
%
\setlength{\columnheight}{104cm}           % value for portrait 3-col mode
\setlength{\columnheighta}{\columnheight-6cm}
\setlength{\columnheightb}{\columnheight+6.5cm}
\setlength{\columnwidth}{0.31\textwidth}   % value for portrait 3-col mode
%
\setlength{\fboxsep}{10mm}
\setlength{\fboxrule}{0mm}
%
%
%
% the page header
%
\hspace*{-1cm}\begin{minipage}[t][][t]{\textwidth-2\fboxsep-8\fboxrule}
\color{invheading}
\begin{center}
{
\VERYHuge \vspace*{-0.6cm}
\textsf{\textbf{
Lab::Measurement -- measurement control with Perl
}}
}
\\[\baselineskip]
\begin{tabular}{ccc}
\parbox{5cm}{\vspace*{-1cm}
\begin{center}
\epsfig{file=images/logo-ur,height=4.5cm}
\end{center}
}
&
\parbox{70cm}{
\begin{center} 
{
\vspace*{-0.5cm}
\huge 
S. Reinhardt$^1$, C. Butschkow$^1$, A. Iankilevitch$^1$, A. Dirnaichner$^1$, 
and A. K. Hüttel$^1$
}
\\[\medskipamount]
\it 
\Large
$^1$Institute for Experimental and Applied
Physics, University of Regensburg, 93040 Regensburg, Germany\\
\end{center}
\vspace*{0.5cm}
}
&
\parbox{5cm}{\vspace*{-1cm}
\begin{center}
\epsfig{file=images/logo-sfb689,height=4.5cm}
\end{center}
}
\end{tabular}
\end{center}
\end{minipage}
\vspace*{9mm}





%
% begin of first row first column
\fcolorbox{white}{background}{
\begin{minipage}[t][\columnheighta-2\fboxsep-2\fboxrule][t]
  {\columnwidth-2\fboxsep-2\fboxrule} \rule{0pt}{0pt}\\
  \begin{minipage}{\textwidth}
%
%

\vspace*{0.9cm}
\heading{Flexible measurement needed?!}

\vspace*{-1cm}
\begin{itemize}
\item
Tired of following your wires in square meters of LabVIEW diagrams?
\item
Tired of clumsy string handling and low-level driver functions
in your looong C program?
\item
Use a text processing language to manage your measurement! Use Perl!
\end{itemize}

\placelisting{srs_read.pl}

\vspace*{2cm}
\vspace*{0.9cm}
\heading{Currently supported hardware}

\vspace*{-3mm}
\placefigure{file=images/hardware}
Hardware driver backends: 
\begin{itemize}
\item
LinuxGPIB and all hardware supported by it
\item
Linux USB-TMC kernel driver
 \item 
NI-VISA (Windows and Linux) and all hardware supported by it
 \item
Oxford Instruments IsoBus
 \item
TCP connection, generic network socket
 \item
RS232 serial port
\end{itemize}

Growing number of High-level drivers (more are {\it very} easy to add):
{\large
\begin{itemize}
\item Various HP/Agilent/Keysight multimeters.
\item DC sources from Yokogawa and Keithley.
\item Lock In Amplifiers from Stanford Research and Signal Recovery.
\item Temperature controllers from Lakeshore and Oxford Instruments.
\item Rohde \& Schwarz RF/Microwave Sources, Spectrum Analyzers, VNAs.
\item and many more...
\end{itemize}
}
\vspace*{1cm}
\heading{Key facts}
\vspace*{-2cm}
\begin{itemize}
 \item
 open source / free software\\[-5cm]
 \item 
 http://www.labmeasurement.de/
 \hfill\epsfig{file=images/qrcode-labmeasurement,width=5cm}
 \item
 license: same as Perl (GPL-1+ or Artistic)
 \item
 releases on CPAN, development on Github
 \item
 contributors and cooperations welcome!
\end{itemize}






%
%
%
\end{minipage}
%
\end{minipage}}
% end of first row first column
%
%
%
%
% begin of first row second column
\fcolorbox{white}{background}{
\begin{minipage}[t][\columnheighta-2\fboxsep-2\fboxrule][t]
  {\columnwidth-2\fboxsep-2\fboxrule} \rule{0pt}{0pt}\\
  \begin{minipage}{\textwidth}
%
%
%
%

\vspace*{0.9cm}

\heading{Real world measurement}
\vspace*{-3mm}
Nested sweeps of gate voltage $V_g$ and bias voltage $V_{sd}$.
Read out multimeter for each point, write values to datafile and create live
color plot:

\placesmalllisting{xpress_2d_example.pl}

\vspace*{2cm}
Script creates output folder:
\vspace*{5mm}
\placesmalllisting{ls.txt}
\vspace*{2cm}
\begin{itemize}
\item archival copy of the measurement script
\item data, tab-separated gnuplot format
\item live plot at end of measurement
\end{itemize}




%
%
%
%
\end{minipage}
\end{minipage}}
% end of second column
%
%
%
%
%
%begin third column
\fcolorbox{white}{background}{
\begin{minipage}[t][\columnheighta-2\fboxsep-2\fboxrule][t]
  {\columnwidth-2\fboxsep-2\fboxrule} \rule{0pt}{0pt}\\
  \begin{minipage}{\textwidth}
%
%
%
%
\vspace*{0.9cm}

\heading{Internal architecture}

\vspace*{-3mm}
\placefigure{file=images/structure}

\begin{itemize}
\item
Modular structure. Easy to extend with new instrument drivers and connection
types.
\item
Abstract IO-layer makes instrument drivers independent of hardware backends.
\end{itemize}
\vspace*{0.9cm}
\heading{Recent additions}
\vspace*{-2cm}
\begin{itemize}
\item New drivers for R\&S FSV Spectrum analyzer and  R\&S ZVA and ZVM vector
  network analyzers, the HP33120 function generator and the Tektronix TDS2024
  oscilloscope.
\item Automated testing of instrument drivers and high-level functionality via
  mock connection objects. Continuous
  integration testing on both Linux (Travis CI) and Windows (Appveyor).
\item New framework for drivers and connections based on \textbf{Moose}, which
  enables 
  modern object oriented programming (OOP) with Perl5:
  \begin{itemize}
  \item Less boilerplate, create classes in a descriptive way. New de-facto
    standard for OOP with Perl5.
  \item Use \textbf{Roles} for safe and efficient sharing of functionality
    between 
    drivers.
  \item Organize SCPI subsystems into Roles, which can be used by multiple
    drivers.
  \item Extensive type system. Convenient validation of user provided
    function parameters.
  \end{itemize}
  
\item \textbf{PDL::Graphics::Gnuplot} gives convenient access to all of
  gnuplot's plot and curve options. Can be used for live plots (qt, x11, wxt)
  and hardcopies (png, jpeg, svg, tikz, pdfcairo, ...).
\end{itemize}

\vspace*{0.9cm}
\heading{Outlook}
\vspace*{-2cm}

\begin{itemize}
\item Implement more drivers, connections and XPRESS functionality with
  Moose.
\item Build portable, lightweight USB-TMC connection based on libusb.
\end{itemize}


\end{minipage}

\end{minipage}}  % end of first row third column
%
%
%

\vspace*{0.5cm}
\begin{center}\color{invheading}\large
\hspace*{-4cm}We gratefully acknowlegde funding by the DFG via the Emmy Noether grant
Hu1808/1-1 ``Carbon nanotubes as electronical and nano-electromechanical hybrid systems in
the quantum limit'',
\\ \hspace*{-4cm}
the collaborative research centre SFB 689 ``Spin phenomena in reduced dimensions'', and
the graduate research school GRK 1570  ``Electronic properties of
carbon based nanostructures''.
\end{center}

\end{document}
